\chapter{Giới thiệu đề
tài}\label{giux1edbi-thiux1ec7u-ux111ux1ec1-tuxe0i}

\section{Đặt vấn đề}\label{ux111ux1eb7t-vux1ea5n-ux111ux1ec1}

Trong bối cảnh bùng nổ của mạng xã hội, đặc biệt là TikTok và Facebook
tại thị trường Việt Nam, hoạt động marketing đóng vai trò sống còn đối
với các doanh nghiệp vừa và nhỏ (SME) ngành Thực phẩm \& Đồ uống (F\&B).
Tuy nhiên, các chủ doanh nghiệp F\&B thường là chuyên gia trong lĩnh vực
ẩm thực nhưng lại thiếu kiến thức, kỹ năng và thời gian để xây dựng,
triển khai một chiến lược marketing số nhất quán và hiệu quả. Họ gặp khó
khăn trong việc sáng tạo nội dung hấp dẫn, quản lý lịch đăng bài đa
kênh, và tương tác kịp thời với khách hàng.

Các giải pháp marketing agency truyền thống quá tốn kém, trong khi các
công cụ quản lý mạng xã hội hiện tại thường yêu cầu người dùng phải tự
tạo nội dung và quản lý workflow thủ công, gây tốn kém thời gian và
không giải quyết được vấn đề thiếu ý tưởng sáng tạo. Sự thiếu nhất quán
trong nội dung và quản lý chiến dịch dẫn đến việc bỏ lỡ cơ hội thu hút
khách hàng mới và giữ chân khách hàng hiện tại.

Trong bối cảnh đó, việc ứng dụng công nghệ AI Agent dựa trên các mô hình
ngôn ngữ lớn (LLM) mở ra một hướng giải quyết đột phá. Một AI Agent
chuyên biệt có thể hoạt động như một "bộ não thứ hai" (second brain), tự
động thu thập cảm hứng từ các nguồn bên ngoài (YouTube, blog), biến
chúng thành các chiến lược và nội dung marketing chuẩn chỉnh, đồng thời
tự động đăng tải lên các nền tảng (Facebook, TikTok).

Vì vậy, việc nghiên cứu và phát triển một AI Agent tự động hóa marketing
chuyên biệt cho SME ngành F\&B là vô cùng cần thiết và có ý nghĩa thực
tiễn. Hệ thống này không chỉ giúp chủ doanh nghiệp tiết kiệm thời gian,
tối ưu hóa nguồn lực mà còn nâng cao hiệu quả chiến dịch marketing, thúc
đẩy doanh thu và tăng cường sự gắn kết của khách hàng đối với thương
hiệu.

\section{Mục tiêu của đề
tài}\label{mux1ee5c-tiuxeau-cux1ee7a-ux111ux1ec1-tuxe0i}

\subsection{Mục đích của đề
tài}\label{mux1ee5c-ux111uxedch-cux1ee7a-ux111ux1ec1-tuxe0i}

Phát triển một AI Agent nhằm hỗ trợ các chủ doanh nghiệp F\&B tự động
hóa quy trình tạo nội dung marketing từ các nguồn cảm hứng thu thập được
trên Internet.

Hệ thống giúp người dùng, những người không chuyên về marketing, chuyển
đổi các nguồn cảm hứng như video, bài viết và hình ảnh thành các nội
dung marketing chuyên nghiệp, bao gồm caption, kịch bản bài đăng và kế
hoạch nội dung có cấu trúc.

Hệ thống cung cấp các nội dung hoàn chỉnh dưới dạng văn bản và hình ảnh
để chủ doanh nghiệp có thể xem xét, chỉnh sửa và tự đăng tải lên các nền
tảng mạng xã hội theo nhu cầu của mình.

Mục tiêu là giúp các chủ doanh nghiệp F\&B tạo ra các nội dung marketing
nhất quán, có chất lượng cao mà không cần phải tự xây dựng ý tưởng hoặc
viết nội dung một cách thủ công.

\subsection{Mục tiêu cụ thể}\label{mux1ee5c-tiuxeau-cux1ee5-thux1ec3}

Xây dựng và chứng minh tính khả thi của hệ thống AI Agent có khả năng:

\begin{itemize}
\item
  Thu thập và xử lý thông tin:

  \begin{itemize}
  \item
    Cho phép người dùng tải lên và lưu trữ các nguồn cảm hứng bao gồm
    văn bản, tài liệu, URL và các nội dung thu thập từ Internet.
  \item
    Sử dụng các kỹ thuật xử lý dữ liệu để trích xuất, phân tích và tổ
    chức thông tin từ các nguồn phi cấu trúc này, tạo thành cơ sở tri
    thức có thể truy xuất.
  \end{itemize}
\item
  Tạo nội dung marketing:

  \begin{itemize}
  \item
    Tự động tổng hợp và chuyển đổi các thông tin đã lưu trữ thành các
    nội dung marketing cụ thể bao gồm caption, kịch bản bài đăng và kế
    hoạch nội dung.
  \item
    Đảm bảo các nội dung được tạo ra có ngữ pháp chính xác, giọng văn
    phù hợp với đặc điểm thương hiệu và phù hợp với thị trường Việt Nam.
  \end{itemize}
\item
  Cung cấp và quản lý nội dung:

  \begin{itemize}
  \item
    Cung cấp các nội dung marketing hoàn chỉnh dưới dạng văn bản có cấu
    trúc, cho phép người dùng xem xét, chỉnh sửa và sử dụng theo nhu
    cầu.
  \item
    Cho phép người dùng quản lý các nội dung đã được tạo ra, bao gồm
    việc xem lại, chỉnh sửa và tổ chức các nội dung theo các chiến dịch
    hoặc thời gian cụ thể.
  \end{itemize}
\item
  Đánh giá và cải thiện:

  \begin{itemize}
  \item
    Đánh giá trải nghiệm người dùng để đảm bảo hệ thống dễ sử dụng đối
    với các chủ doanh nghiệp F\&B không có chuyên môn về công nghệ.
  \item
    Thiết lập cơ chế thu thập phản hồi từ người dùng về chất lượng nội
    dung được tạo ra, nhằm cải thiện khả năng sinh nội dung của hệ
    thống.
  \end{itemize}
\end{itemize}

\subsection{Thách thức và phương pháp tiếp
cận}\label{thuxe1ch-thux1ee9c-vuxe0-phux1b0ux1a1ng-phuxe1p-tiux1ebfp-cux1eadn}

\subsubsection{Thách thức}\label{thuxe1ch-thux1ee9c}

\begin{itemize}
\item
  Thiếu kỹ năng và thời gian của chủ quán F\&B:

  \begin{itemize}
  \item
    Chủ quán là người không chuyên về marketing và công nghệ, yêu cầu
    giao diện cực kỳ đơn giản, "zero-config" (không cần cấu hình phức
    tạp).
  \item
    Họ rất bận rộn, cần một giải pháp tự động hóa cao độ, giảm thiểu tối
    đa sự can thiệp thủ công hàng ngày.
  \end{itemize}
\item
  Chất lượng nội dung và tính xác thực (Authenticity):

  \begin{itemize}
  \item
    Nội dung marketing F\&B đòi hỏi sự chân thực (authentic footage),
    trong khi AI tạo sinh hình ảnh/video thường kém chân thực hoặc tốn
    kém.
  \item
    Thách thức là làm sao nâng cao chất lượng nội dung từ nguồn đầu vào
    không chuyên nghiệp của chủ quán (ảnh chụp bằng điện thoại thiếu
    sáng) và đảm bảo giọng văn (tone of voice) phù hợp, không bị "máy
    móc".
  \end{itemize}
\item
  Bắt kịp xu hướng và tính địa phương hóa (Localization):

  \begin{itemize}
  \item
    Xu hướng F\&B và social media (TikTok trend, hashtag) thay đổi liên
    tục, đặc biệt tại thị trường Việt Nam.
  \item
    Hệ thống cần liên tục cập nhật dữ liệu xu hướng để đề xuất nội dung
    phù hợp, tránh tạo ra nội dung lỗi thời hoặc không liên quan đến văn
    hóa địa phương.
  \end{itemize}
\end{itemize}

\subsubsection{Phương pháp tiếp
cận}\label{phux1b0ux1a1ng-phuxe1p-tiux1ebfp-cux1eadn}

\begin{itemize}
\item
  Sử dụng AI Agent dựa trên LLM để xây dựng chiến lược và sáng tạo nội
  dung:

  \begin{itemize}
  \item
    LLM đóng vai trò cốt lõi trong việc phân tích nội dung cảm hứng đã
    lưu (YouTube, blog), hiểu ngữ cảnh F\&B và chuyển hóa thành các ý
    tưởng bài viết/kịch bản video có cấu trúc.
  \item
    Áp dụng kỹ thuật Prompt Engineering để định hình các "personas"
    (giọng văn) marketing chuyên nghiệp, giúp chủ quán dễ dàng lựa chọn
    mà không cần cấu hình phức tạp.
  \end{itemize}
\item
  Thiết kế luồng công việc tinh gọn, lấy người dùng làm trung tâm:

  \begin{itemize}
  \item
    Xây dựng giao diện người dùng đơn giản, tập trung vào việc cho phép
    người dùng tải lên và quản lý các nguồn cảm hứng, sau đó xem xét,
    chỉnh sửa và sử dụng các nội dung đã được hệ thống tạo ra.
  \item
    Tối ưu hóa quy trình làm việc để giảm thiểu các thao tác thủ công,
    cho phép người dùng dễ dàng chuyển đổi từ việc thu thập nguồn cảm
    hứng sang việc nhận được các nội dung marketing hoàn chỉnh.
  \end{itemize}
\item
  Xây dựng và quản lý cơ sở dữ liệu nội dung:

  \begin{itemize}
  \item
    Thiết kế hệ thống lưu trữ dựa trên cơ sở dữ liệu vector (Vector
    Database) để tổ chức và quản lý hiệu quả các nguồn cảm hứng và nội
    dung đã được xử lý.
  \item
    Cho phép người dùng quản lý tập hợp các tài liệu và nội dung đã được
    tạo ra, bao gồm việc thêm mới, xem lại và tổ chức các nội dung theo
    các chiến dịch hoặc thời gian cụ thể.
  \end{itemize}
\item
  Triển khai và đánh giá prototype:

  \begin{itemize}
  \item
    Xây dựng và triển khai một prototype hoàn chỉnh bao gồm các chức
    năng thu thập, xử lý, phân tích nguồn cảm hứng và tạo nội dung
    marketing.
  \item
    Tiến hành thử nghiệm prototype với người dùng thực tế, chủ yếu là
    các chủ doanh nghiệp F\&B, để đánh giá hiệu quả của hệ thống trong
    việc tạo ra các nội dung marketing có chất lượng và phù hợp với nhu
    cầu của họ.
  \item
    Thu thập phản hồi từ người dùng để đánh giá tính khả dụng của hệ
    thống và hiệu quả của các nội dung được tạo ra, từ đó cải thiện chất
    lượng và trải nghiệm người dùng.
  \end{itemize}
\end{itemize}

\section{Bố cục báo cáo}\label{bux1ed1-cux1ee5c-buxe1o-cuxe1o}

Báo cáo được chia thành 5 chương chính:

\begin{itemize}
\item
  \textbf{Chương 1 Giới thiệu đề tài:} Trình bày lý do chọn đề tài, bối
  cảnh, thách thức và mục tiêu nghiên cứu.
\item
  \textbf{Chương 2 Tìm hiểu và nghiên cứu thị trường:} Phân tích thực
  trạng SME, nhu cầu quản lý tri thức và truy xuất thông tin, khảo sát
  các giải pháp hiện có.
\item
  \textbf{Chương 3 Kiến thức nền tảng và công nghệ sử dụng:} Tổng quan
  về các công nghệ liên quan như NLP, LLM, RAG và các kiến thức nền tảng
  để thiết kế AI Agent.
\item
  \textbf{Chương 4 Phân tích và thiết kế hệ thống:} Phân tích các yêu
  cầu chức năng, yêu cầu phi chức năng, yêu cầu giao diện và thiết kế hệ
  thống tổng quát
\item
  \textbf{Chương 5 Kết luận và định hướng:} Tổng kết kết quả nghiên cứu,
  những đóng góp chính của đề tài và định hướng phát triển trong tương
  lai.
\end{itemize}
